\section{本章小结}
\label{sec:nspace-summary}

章首的问题是：比较两个函数时用哪一种误差，以及误差趋零的过程能否在空间内找到终点。范数回答前一个问题。它满足正定性、齐次性和三角不等式，诱导出距离、收敛和柯西列；最大绝对值看最坏位置，积分看总体累计，两者在$C[0,1]$上只有单向的控制$\lVert f\rVert_1\leq\lVert f\rVert_\infty$。有限维中任意两个范数等价，无限维中则不然：借助第~\ref{chap:functions-as-vectors}章两两远离的单位列和一组代数基，可以在任何无限维赋范空间上造出不等价的范数，$x^n$则在$C[0,1]$上给出具体例子。所以无限维中的收敛必须连同范数一起陈述。

完备性回答后一个问题。Banach空间中每个柯西列都有极限；Banach空间的线性子空间完备当且仅当它是闭集，多项式在最大误差下就不闭。紧集上的连续函数在一致范数下完备，证明依靠一致柯西列的极限与一致极限的连续性；非紧定义域上改用有界连续函数，只需补上极限的有界性。同一个$C[0,1]$在积分范数下不完备，逼近阶跃函数的折线列没有连续的极限。

$L^p$与$L^\infty$为积分误差下的极限提供了去处。它们的元素是几乎处处相等的可测函数组成的等价类，因而一般没有点值；计数测度下得到数列空间$\ell^p$与$\ell^\infty$。$L^p$的三角不等式来自Young、Hölder与Minkowski不等式，完备性的证明先抽相邻误差可求和的子列，再用Fatou引理回到积分。$L^\infty$用本质上确界忽略零测集，Hölder不等式在端点$p=1$、$q=\infty$同样成立，完备性来自把可数个例外集合并成一个零测集后的一致收敛。

各种收敛方式控制的误差不同。一致收敛推出$L^p$收敛需要有限测度；逐点收敛管不住窄尖峰，$L^p$收敛推不出逐点收敛，打字机序列是例子；但$L^p$收敛与依测度收敛的序列都有几乎处处收敛的子列。依测度收敛只控制大误差出现的范围，在有限测度下配合一致可积性才推出$L^1$收敛，共同可积上界是一致可积的充分条件，由此回到支配收敛定理。

本章只在单个空间内比较误差。内积带来的正交结构与最佳逼近在第~\ref{chap:inner-product-spaces}章研究；把一个空间放进另一个空间时误差怎样传递，即包含映射的有界性，在第~\ref{chap:linear-operators}章研究；用简单函数、连续函数逼近$L^p$中的元素，在第~\ref{chap:function-approximation}章研究。
